\documentclass[11pt]{article}

\usepackage[margin=1in]{geometry}
\usepackage{amsmath}
\usepackage{booktabs}
\usepackage{array}
\usepackage{hyperref}

\title{White by Construction: Two Crypto Stat-Arbitrage Screens\\
That Pass Their Own Placebos}
\author{Jack Lutz \and Sammy Adham \and Frank Kronewitter}
\date{\today}

\begin{document}

\maketitle

\begin{abstract}
We set out to build a profitable statistical-arbitrage strategy on hourly cryptocurrency
pair spreads. We report mostly what does not work, and why that is the more useful result.
Two screens that appear to find strong, tradable structure are statistical artifacts. A
Kalman dynamic hedge ratio appears to produce stationary, mean-reverting spreads for nearly
the entire liquid universe, but it runs a stationarity test on a filter's own innovations,
which are white by construction; placebos that cannot be cointegrated pass the same screen at
the same near-$100\%$ rate, while a clean test finds genuine out-of-sample cointegration in
only a few percent of pairs. A rolling $z$-score appears to generate a large reversion edge,
but it reverts mechanically even on a pure random walk. Both artifacts hold at hourly,
four-hour, and daily frequencies. Once we benchmark against placebos, the structure that
survives is either priced too fast to trade or too marginal to deploy: order-flow information
in the Level-2 book decays within seconds, reversion speed is a selectable property of pairs
(Spearman $\rho = 0.46$) but selectability is not profit, and the one market-neutral
mean-reversion book that clears honest accounting posts a HAC- and venue-adjusted monthly
Sharpe near $1.0$, and only as a patient, never-stop, deep-drawdown hold. We read the result
as a market-efficiency statement about crypto pair spreads, and as a cautionary standard: a
cointegration or $z$-score screen should not be trusted until it has been shown to fail on a
matched placebo.
\end{abstract}

% =====================================================================
\section{Introduction}
% TODO (write after Sections 3--5 lock): open on the contribution; state the original
% profit question and the changed conclusion in paragraph one; anchor to the proposal's
% market-efficiency framing; list contributions (two artifacts, one efficiency null, one
% bounded effect). Use the Jensen-Kelly-Pedersen hinge sentence.

% =====================================================================
\section{Data and market setting}
% TODO: venues (Binance.US hourly klines for the spread work; Coinbase hourly for the
% cross-exchange check; Tardis Level-2 book + trades, Binance spot, 2024, top-50, for the
% microstructure work), frequencies, costs/slippage stated up front, universe construction,
% and the inference discipline (HAC everywhere, strict walk-forward, a placebo for every
% positive claim).

% =====================================================================
\section{Two screening artifacts}
\label{sec:artifacts}

Two screens drove our early results, and both are artifacts. The first, a Kalman dynamic
hedge ratio, appeared to find out-of-sample cointegration in almost every liquid pair. The
second, a rolling $z$-score, appeared to find strong mean reversion. We show that each passes
a placebo it should fail, and that neither carries the information it seems to. This is the
paper's main methodological contribution. Both screens are in common use, and a careful
analyst is easily fooled by either one.

\subsection{The Kalman dynamic-hedge cointegration screen}
\label{sec:kalman}

A static hedge ratio fits one $\beta$ for the whole sample. A Kalman dynamic hedge lets
$\beta$ drift as a random walk and updates it each step from the data. We fit the process and
observation noise by maximum likelihood on the training window, froze them, and rolled the
filter forward on the test window. The test-period innovations, the filter's one-step
prediction errors, became the spread, and we ran an augmented Dickey--Fuller (ADF) test on
them. Under this screen $100\%$ of the top-50 pairs rejected the unit root out of sample at
$p<0.05$, and $96.4\%$ at $p<0.001$. A static-OLS spread over the same window passed only
about $5\%$ of pairs. The contrast looked decisive.

It is decisive about the filter, not the pair. A Kalman filter with non-zero process noise
adapts its hedge ratio to track whatever relationship is present, so its innovations are white
by construction. An ADF test on a white series rejects the unit root almost always, whether or
not the two assets are cointegrated. Testing the stationarity of a filter's own innovations is
circular. An earlier draft of this project defended the result on the grounds that the
parameters were fit on training data only and never refit on test. That defense misses the
mechanism. The whitening holds out of sample too, because the filter keeps adapting on the
test window from its frozen dynamics.

A placebo settles it. We ran the identical screen on series that cannot be cointegrated:
(a) independent simulated random walks, (b) each coin paired with a phase-randomized surrogate
of another coin, and (c) real pairs with one leg block-shuffled. All three pass at the same
near-$100\%$ rate as the real pairs (Table~\ref{tab:kalman-placebo}). A clean, non-circular
test tells the opposite story. Engle--Granger~\citep{engle1987} on the test window, and a
static-OLS ADF, put genuine out-of-sample cointegration at the null floor, roughly $2$--$5\%$
of pairs on the matched window. Essentially none of the universe is cointegrated out of
sample. The screen's $100\%$ is noise dressed as signal, and the placebo result is, to our
knowledge, the cleanest way to see it~\citep{bailey2014}.

\begin{table}[h]
\centering
\caption{The Kalman screen passes placebos that cannot be cointegrated at the same rate as
real pairs. Out-of-sample ADF on the Kalman innovations, top-50 USDT universe, $t_0 =$
2024-01-01, 90-day train / 30-day test.}
\label{tab:kalman-placebo}
\begin{tabular}{lrr}
\toprule
Series tested & ADF $p<0.05$ & ADF $p<0.001$ \\
\midrule
Real top-50 pairs & $100.0\%$ & $96.4\%$ \\
(a) Independent random walks & $100.0\%$ & $98.3\%$ \\
(b) Coin vs.\ phase-randomized surrogate & $100.0\%$ & $95.0\%$ \\
(c) Real pair, one leg block-shuffled & $100.0\%$ & $95.0\%$ \\
\midrule
Clean Engle--Granger (test window) & $2$--$3\%$ & --- \\
\bottomrule
\end{tabular}
\end{table}

The artifact does not depend on the bar. We reran the screen at four-hour and daily
frequencies (Table~\ref{tab:kalman-freq}). Real pairs and random-walk placebos pass together
at every cadence: $100\%$ versus $100\%$ at hourly and four-hour, and $70\%$ versus $68\%$ at
daily ($p<0.001$). The clean Engle--Granger rate behaves differently. It is a small,
window-dependent fraction that rises over long hourly windows, where crypto majors share a
common trend, and falls to the $5\%$ null floor at daily. It never approaches the screen's
$100\%$. The Kalman screen is uninformative about cointegration at every frequency we tested.

\begin{table}[h]
\centering
\caption{Frequency robustness of the Kalman artifact. Real pairs and random-walk placebos pass
the screen together at every cadence; the clean Engle--Granger benchmark does not. Top-50
universe, 60 pairs per cell, recent multi-year window.}
\label{tab:kalman-freq}
\begin{tabular}{lrrrrr}
\toprule
& \multicolumn{2}{c}{Real pairs} & \multicolumn{2}{c}{RW placebo} & Clean EG \\
\cmidrule(lr){2-3}\cmidrule(lr){4-5}\cmidrule(lr){6-6}
Frequency & $p<0.05$ & $p<0.001$ & $p<0.05$ & $p<0.001$ & $p<0.05$ \\
\midrule
Hourly & $100.0\%$ & $100.0\%$ & $100.0\%$ & $100.0\%$ & $25.0\%$ \\
4-hour & $100.0\%$ & $100.0\%$ & $100.0\%$ & $100.0\%$ & $11.7\%$ \\
Daily  & $100.0\%$ & $70.0\%$  & $100.0\%$ & $68.3\%$  & $5.0\%$ \\
\bottomrule
\end{tabular}
\end{table}

\subsection{The rolling $z$-score}
\label{sec:rollingz}

The second screen is a rolling $z$-score, the standard entry signal for pair trades. For a
spread $s_t$ it subtracts a trailing mean and divides by a trailing standard deviation,
\[
z_t = \frac{s_t - \bar{s}_{t-1}}{\sigma_{t-1}},
\]
and it trades when $|z_t|$ is large, betting the spread returns to its mean. On our spreads it
produced a large reversion edge.

The edge is mechanical. Subtracting a trailing mean from any series, including a pure random
walk, manufactures mean reversion, because the trailing mean chases the level and the deviation
from it cannot wander as freely as the raw series. A random walk has no mean to revert to, yet
a rolling $z$-score on one still reverts.

We measured the mechanical floor directly. On pure random walks, the same rolling-$z$ rule
produces strong positive reversion at every frequency we tested: $+0.98z$ per event at hourly,
$+1.61z$ at four-hour, and $+1.51z$ at daily, each far outside its sampling noise
(Table~\ref{tab:rollingz}). As a trading rule the floor is large. A rolling-$z$ strategy that
posts a Sharpe near $4.75$ on our data still posts about $2.2$ on random pairs and $2.4$ on
phase-randomized noise. Most of the apparent edge is the floor. For this reason every reversion
number elsewhere in this paper is reported net of a per-pair, per-window random-walk floor.

\begin{table}[h]
\centering
\caption{The rolling $z$-score manufactures mean reversion on pure random walks, at every
frequency. Mean signed reversion per $|z|>2$ event, 200 random-walk surrogates per cell.}
\label{tab:rollingz}
\begin{tabular}{lrr}
\toprule
Frequency & Mean event reversion ($z$) & s.d. \\
\midrule
Hourly & $+0.98$ & $0.26$ \\
4-hour & $+1.61$ & $0.39$ \\
Daily  & $+1.51$ & $0.45$ \\
\bottomrule
\end{tabular}
\end{table}

Both screens share a failure mode. They manufacture the appearance of structure that a matched
placebo reproduces. The lesson is narrow and practical. A cointegration screen built on a
filter's own innovations must be checked against a white-noise placebo before it is believed,
and a rolling-$z$ reversion claim must be checked against a random-walk null. Neither check is
standard practice, and without them both screens report a signal that is not there.

% =====================================================================
\section{What is real but not tradable}
% TODO: reversion selectability (Spearman rho=0.46, positive in all 9 splits, clean placebos;
% selectability is not profit) and the microstructure efficiency null (order-flow information
% with a seconds-scale half-life; the execution null). One exhibit: order-flow half-life decay.

% =====================================================================
\section{The marginal survivor}
% TODO: the never-stop mean-reversion book at a HAC- and venue-honest monthly Sharpe ~1.0,
% market-neutral, selection-sensitive, patient-only. Lead with the bound, not the point
% estimate. Period-partitioned table, not a single headline Sharpe. Guardrail: do not
% over-claim.

% =====================================================================
\section{Limitations}
% TODO: dedicated threats-to-validity. Point-in-time dead-coin universe still open;
% selection-bias mitigated not closed (no single pre-registered config evaluated once);
% single venue / quote; microstructure is 2024-only and <=4 symbols for event-level work.

% =====================================================================
\section{Conclusion}
% TODO: modest, market-efficiency framed. End on a standard the field should meet (any
% cointegration or z-score screen must pass a matched placebo before it is believed). Extend,
% do not restate.

% =====================================================================
\begin{thebibliography}{9}

\bibitem{engle1987}
Engle, R.~F., and Granger, C.~W.~J. (1987).
Co-integration and error correction: representation, estimation, and testing.
\textit{Econometrica}, 55(2), 251--276.

\bibitem{gatev2006}
Gatev, E., Goetzmann, W.~N., and Rouwenhorst, K.~G. (2006).
Pairs trading: performance of a relative-value arbitrage rule.
\textit{Review of Financial Studies}, 19(3), 797--827.

\bibitem{bailey2014}
Bailey, D.~H., and L\'opez de Prado, M. (2014).
The deflated Sharpe ratio: correcting for selection bias, backtest overfitting, and
non-normality.
\textit{Journal of Portfolio Management}, 40(5), 94--107.

\bibitem{harvey2016}
Harvey, C.~R., Liu, Y., and Zhu, H. (2016).
$\ldots$ and the cross-section of expected returns.
\textit{Review of Financial Studies}, 29(1), 5--68.

\end{thebibliography}

\end{document}
